\section{Predicción del parámetro de área}

En el modelo clásico de Fay--Herriot, el objetivo principal no es estimar únicamente el vector de parámetros fijos $\beta$, sino \textbf{predecir el verdadero parámetro de interés del área} $i$, denotado por $\theta_i$.

Recordemos que el modelo se especifica en dos niveles:
\[
y_i=\theta_i+e_i, \qquad e_i\sim N(0,D_i),
\]
\[
\theta_i=x_i^{\top}\beta+u_i, \qquad u_i\sim N(0,A),
\]
con independencia entre $e_i$ y $u_i$.

Sustituyendo la segunda ecuación en la primera, obtenemos la forma combinada del modelo:
\[
y_i=x_i^{\top}\beta+u_i+e_i.
\]

Por lo tanto, el problema de predicción consiste en construir un predictor adecuado de $\theta_i$ a partir de la información observada $y_i$.

\subsection{¿Qué necesitamos para la predicción?}

Para derivar el predictor de $\theta_i$, primero necesitamos calcular:
\[
E(\theta_i), \qquad E(y_i), \qquad \operatorname{Cov}(\theta_i,y_i), \qquad \operatorname{Var}(y_i).
\]

\subsubsection*{Esperanza de $\theta_i$}

Como
\[
\theta_i=x_i^{\top}\beta+u_i,
\]
y $x_i^{\top}\beta$ es una cantidad fija (no aleatoria), entonces
\[
E(\theta_i)=E(x_i^{\top}\beta+u_i)=x_i^{\top}\beta+E(u_i).
\]
Dado que
\[
u_i\sim N(0,A),
\]
se tiene $E(u_i)=0$, y por tanto
\[
E(\theta_i)=x_i^{\top}\beta.
\]

\subsubsection*{Esperanza de $y_i$}

Usando la forma combinada
\[
y_i=x_i^{\top}\beta+u_i+e_i,
\]
tenemos
\[
E(y_i)=E(x_i^{\top}\beta+u_i+e_i)=x_i^{\top}\beta+E(u_i)+E(e_i).
\]
Como
\[
E(u_i)=0, \qquad E(e_i)=0,
\]
resulta que
\[
E(y_i)=x_i^{\top}\beta.
\]

\subsubsection*{Covarianza entre $\theta_i$ y $y_i$}

Ahora calculemos
\[
\operatorname{Cov}(\theta_i,y_i).
\]
Sustituyendo las expresiones de $\theta_i$ y $y_i$,
\[
\
\operatorname{Cov}(x_i^{\top}\beta+u_i,\;x_i^{\top}\beta+u_i+e_i)
\]

\[
\
\operatorname{Cov}(x_i^{\top}\beta,x_i^{\top}\beta)
+\operatorname{Cov}(x_i^{\top}\beta,u_i)
+\operatorname{Cov}(x_i^{\top}\beta,e_i)
+\operatorname{Cov}(u_i,x_i^{\top}\beta)
+\operatorname{Cov}(u_i,u_i)
+\operatorname{Cov}(u_i,e_i).
\]



Usamos las siguientes propiedades de la covarianza:
\begin{itemize}
    \item la covarianza de una constante con cualquier variable aleatoria es cero;
    \item la covarianza es bilineal;
    \item si dos variables son independientes, su covarianza es cero.
\end{itemize}

Entonces,
\[
\operatorname{Cov}(x_i^{\top}\beta+u_i,\;x_i^{\top}\beta+u_i+e_i)
=
\operatorname{Cov}(u_i,u_i)+\operatorname{Cov}(u_i,e_i),
\]
porque todos los términos que involucran la constante $x_i^{\top}\beta$ desaparecen.

Además, como $u_i$ y $e_i$ son independientes,
\[
\operatorname{Cov}(u_i,e_i)=0.
\]
Y como
\[
\operatorname{Cov}(u_i,u_i)=\operatorname{Var}(u_i)=A,
\]
obtenemos
\[
\operatorname{Cov}(\theta_i,y_i)=A.
\]

\subsubsection*{Varianza de $y_i$}

Partimos de
\[
y_i=x_i^{\top}\beta+u_i+e_i.
\]
Entonces
\[
\operatorname{Var}(y_i)=\operatorname{Var}(x_i^{\top}\beta+u_i+e_i).
\]

Usamos ahora la fórmula general
\[
\operatorname{Var}(X+Y)=\operatorname{Var}(X)+\operatorname{Var}(Y)+2\operatorname{Cov}(X,Y),
\]
y su extensión a más de dos términos. Como $x_i^{\top}\beta$ es constante, su varianza es cero. Por tanto,
\[
\operatorname{Var}(y_i)=\operatorname{Var}(u_i)+\operatorname{Var}(e_i)+2\operatorname{Cov}(u_i,e_i).
\]

Por independencia,
\[
\operatorname{Cov}(u_i,e_i)=0.
\]
Además,
\[
\operatorname{Var}(u_i)=A, \qquad \operatorname{Var}(e_i)=D_i.
\]
Así,
\[
\operatorname{Var}(y_i)=A+D_i.
\]

En resumen, hemos obtenido las cuatro cantidades clave:
\[
E(\theta_i)=x_i^{\top}\beta,
\qquad
E(y_i)=x_i^{\top}\beta,
\qquad
\operatorname{Cov}(\theta_i,y_i)=A,
\qquad
\operatorname{Var}(y_i)=A+D_i.
\]

\subsection{Obtención de la esperanza condicional \(E(\theta_i\mid y_i)\)}

Bajo el modelo normal, el vector aleatorio
\[
\begin{pmatrix}
\theta_i\\
y_i
\end{pmatrix}
\]
tiene distribución normal bivariada. Para variables conjuntamente normales, la esperanza condicional viene dada por la fórmula general:
\[
E(U\mid V)
=
E(U)+\operatorname{Cov}(U,V)\operatorname{Var}(V)^{-1}\bigl(V-E(V)\bigr).
\]

Aplicando esta fórmula con
\[
U=\theta_i,
\qquad
V=y_i,
\]
obtenemos
\[
E(\theta_i\mid y_i)
=
E(\theta_i)
+
\operatorname{Cov}(\theta_i,y_i)\operatorname{Var}(y_i)^{-1}\bigl(y_i-E(y_i)\bigr).
\]

Ahora sustituimos, paso a paso, los resultados ya calculados:
\[
E(\theta_i\mid y_i)
=
x_i^{\top}\beta
+
A(A+D_i)^{-1}\bigl(y_i-x_i^{\top}\beta\bigr).
\]

Como en este caso todo es escalar, podemos escribir
\[
(A+D_i)^{-1}=\frac{1}{A+D_i},
\]
de modo que
\[
E(\theta_i\mid y_i)
=
x_i^{\top}\beta
+
\frac{A}{A+D_i}\bigl(y_i-x_i^{\top}\beta\bigr).
\]

Distribuyendo el factor,
\[
E(\theta_i\mid y_i)
=
x_i^{\top}\beta
+
\frac{A}{A+D_i}y_i
-
\frac{A}{A+D_i}x_i^{\top}\beta.
\]

Agrupando los términos que multiplican a $x_i^{\top}\beta$,
\[
E(\theta_i\mid y_i)
=
\frac{A}{A+D_i}y_i
+
\left(1-\frac{A}{A+D_i}\right)x_i^{\top}\beta.
\]

Definiendo
\[
\gamma_i=\frac{A}{A+D_i},
\]
obtenemos
\[
\tilde{\theta}_i
=
E(\theta_i\mid y_i)
=
\gamma_i y_i + (1-\gamma_i)x_i^{\top}\beta.
\]

Esta expresión muestra que el predictor de $\theta_i$ es una combinación ponderada entre:
\begin{itemize}
    \item el estimador directo $y_i$;
    \item la parte explicada por el modelo $x_i^{\top}\beta$.
\end{itemize}

\subsection{Interpretación del factor de encogimiento}

El término
\[
\gamma_i=\frac{A}{A+D_i}
\]
se denomina \textbf{factor de encogimiento} o \textit{shrinkage factor}.

Su interpretación es inmediata:

\paragraph{Si $D_i$ es pequeño.}
Cuando la varianza muestral del estimador directo es baja, el estimador directo es preciso. En ese caso,
\[
\gamma_i \approx 1,
\]
y por tanto
\[
\tilde{\theta}_i \approx y_i.
\]
Es decir, la predicción se apoya principalmente en la información directa de la encuesta.

\paragraph{Si $D_i$ es grande.}
Cuando la varianza muestral es alta, el estimador directo es muy ruidoso. En ese caso,
\[
\gamma_i \approx 0,
\]
y entonces
\[
\tilde{\theta}_i \approx x_i^{\top}\beta.
\]
Es decir, la predicción se apoya más en la estructura del modelo.

Así, el estimador de Fay--Herriot ajusta automáticamente el peso dado al estimador directo según la precisión muestral del área.

\subsection{¿Qué es el BLUP?}

Hasta ahora hemos obtenido
\[
E(\theta_i\mid y_i)
=
x_i^{\top}\beta
+
\frac{A}{A+D_i}\bigl(y_i-x_i^{\top}\beta\bigr),
\]
suponiendo conocidos $\beta$ y $A$.

Sin embargo, en la práctica $\beta$ casi nunca es conocido. Por ello, debemos estimarlo a partir de los datos. Si $A$ se considera conocido, entonces la matriz de covarianza marginal de $y$ es
\[
V=AI_m+D,
\]
y el estimador natural de $\beta$ es el estimador de mínimos cuadrados generalizados:
\[
\hat{\beta}(A)=
(X^{\top}V^{-1}X)^{-1}X^{\top}V^{-1}y.
\]

Al reemplazar $\beta$ por $\hat{\beta}(A)$ en el predictor anterior, obtenemos
\[
\tilde{\theta}_i(A)
=
x_i^{\top}\hat{\beta}(A)
+
\gamma_i\bigl(y_i-x_i^{\top}\hat{\beta}(A)\bigr),
\qquad
\gamma_i=\frac{A}{A+D_i}.
\]

Equivalentemente,
\[
\tilde{\theta}_i(A)
=
\gamma_i y_i + (1-\gamma_i)x_i^{\top}\hat{\beta}(A).
\]

A este predictor se le llama \textbf{BLUP} (\textit{Best Linear Unbiased Predictor}).

\paragraph{¿Por qué se llama BLUP?}
Se llama así porque, dentro de la clase de predictores:
\begin{itemize}
    \item lineales en los datos observados,
    \item insesgados bajo el modelo,
\end{itemize}
es el que minimiza el error cuadrático medio de predicción cuando $A$ es conocido.

En otras palabras, el BLUP es la versión ``óptima'' del predictor lineal insesgado bajo el modelo de Fay--Herriot con componente de varianza conocida.

\subsection{¿Qué es el EBLUP?}

En aplicaciones reales, el problema es que $A$ tampoco es conocido. Por tanto, no podemos calcular exactamente ni
\[
\gamma_i=\frac{A}{A+D_i},
\]
ni
\[
\hat{\beta}(A).
\]

La solución consiste en estimar $A$ a partir de los datos, obteniendo un valor $\hat A$ mediante procedimientos como máxima verosimilitud (ML) o máxima verosimilitud restringida (REML).

Una vez obtenido $\hat A$, reemplazamos $A$ por $\hat A$ en todas las expresiones anteriores:
\[
\hat{\gamma}_i=\frac{\hat A}{\hat A+D_i},
\]
\[
\hat{\beta}(\hat A)=
(X^{\top}\hat V^{-1}X)^{-1}X^{\top}\hat V^{-1}y,
\qquad
\hat V=\hat A I_m+D.
\]

Entonces el predictor queda
\[
\hat{\theta}_i^{EBLUP}
=
x_i^{\top}\hat{\beta}(\hat A)
+
\hat{\gamma}_i\bigl(y_i-x_i^{\top}\hat{\beta}(\hat A)\bigr).
\]

Equivalentemente,
\[
\hat{\theta}_i^{EBLUP}
=
\hat{\gamma}_i y_i + (1-\hat{\gamma}_i)x_i^{\top}\hat{\beta}(\hat A).
\]

A este predictor se le denomina \textbf{EBLUP} (\textit{Empirical Best Linear Unbiased Predictor}).

\paragraph{¿En qué se diferencia del BLUP?}
La diferencia esencial es la siguiente:
\begin{itemize}
    \item el \textbf{BLUP} asume conocido el componente de varianza $A$;
    \item el \textbf{EBLUP} reemplaza ese valor desconocido por un estimador $\hat A$.
\end{itemize}

Por eso se le llama \textit{empirical}: porque ya no usa el valor teórico verdadero de $A$, sino una estimación obtenida empíricamente a partir de los datos.

\subsection{Conclusiones}

La predicción del parámetro de área en Fay--Herriot no consiste ni en copiar el estimador directo ni en sustituirlo por una simple regresión. Lo que hace el modelo es \textbf{combinar de manera óptima la información directa y la información auxiliar}.

Así, el predictor final:
\begin{itemize}
    \item conserva la información específica del área mediante $y_i$;
    \item estabiliza la estimación mediante $x_i^{\top}\beta$;
    \item y ajusta automáticamente el peso de cada componente según la precisión muestral del área.
\end{itemize}

Esta es la razón por la cual el modelo de Fay--Herriot produce estimaciones más estables que las estimaciones directas cuando se trabaja con áreas pequeñas.